\section{Limits of the Drop Design}\label{sec:reach}

A benchmark can be missing from a model card because the evaluation was irrelevant to the model, because it was never run, or because it was run and withheld. Only the third is disclosure in the sense of \citet{grossman1981}, \citet{milgrom1981} and \citet{dye1985}, and it is the only one a reader is entitled to price. The within-family drop is built to isolate it. Consider a provider that reports benchmark $b$ for release $t$ and omits $b$ for release $t+1$, where an independent score for release $t+1$ on $b$ exists in the panel. Reporting $b$ once establishes that the provider considered the benchmark relevant and possessed the harness to run it, retiring the first two explanations by the provider's own earlier behaviour rather than by assumption, and the independent score supplies the counterfactual that clinical outcome-reporting work takes from registries \citep{chan2004,dwan2013,kirkham2010}.

The ceiling on that design is small. The panel contains 144 adjacent same-family release pairs, covering 711 shared eligible benchmarks and 15 organisations. Requiring both releases in a pair to clear the eight-benchmark analysis sample, the threshold every other estimate in this paper respects, leaves 37 pairs, 482 shared benchmarks and 8 organisations. The restricted count binds, since a pair resting on a handful of benchmarks cannot support a percentile comparison, and eight organisations is also the effective cluster count, which puts the design in the regime where \autoref{sec:data-inference} reports no analytic standard error.

Precision at that ceiling follows from the null calibration of a drop gap. Under the null of no selective withholding, the sampling standard deviation of the gap is 19.4 percentile points when one benchmark is dropped, 14.4 when two are, and 12.4 when three are. At conventional two-sided size a single-drop event must therefore move the gap by roughly 38 percentile points before it separates from noise. The contamination the uncorrected comparison generates is itself only $+12.23$ points, so a disclosure gap of that size would sit well below the threshold, and no single drop event can be read alone.

What the coded sample supports today is less than that. Disclosure coding covers 48 cells drawn from two Anthropic releases, Claude Opus 4.6 and 4.7, of which 15 are reported and 33 are omitted, distributed across ORBIT categories as 33 in H, 12 in A, 2 in B and 1 in C. No cell falls in category E, the classification's drop category, and the drop estimator accordingly has $n = 0$. This is a case series that demonstrates the instrument and fixes its coding conventions, and it supports no inference of any kind. The boundary is a property of the design space, in that the drop design is the comparison the problem calls for, it is available at most 37 times, and its detection threshold is now known in advance to whoever codes it next.
